\section{Модели}

Рассмотрим базовую задачу организации синхронного движения самолетов по
наземному пространству аэропорта от ворот до выходов - до условной взлётно посадочной полосы \cite{tien2019deep}.
Проезжая часть представлена в виде матрицы с четырьмя воротами (на левом краю матрицы).
 и двумя выходами (на правом краю матрицы).

Всем самолетам-агентам разрешено движение только через указанную общую проезжую часть,
отмеченную зелеными ячейками на сетке и соответствующим положительным значением в матрице (=1) (см. рис \ref{fig:state_map}).
Если полет наступит на запрещенную клетку - фиолетовую ячейку карты с нулевым значением в матрице -
он получит штраф $\phi$.
Любые пересечения между траекториями полетов, например, положение любых двух самолетов
в одной и той же ячейке или на пересечении унитарных движений на карте запрещены
и наказываются штрафом $\beta$.
Количество выходов намеренно выбрано меньшим, чем количество ворот, 
чтобы обеспечить тесное взаимодействие между агентами.

\begin{figure}[h]
    \centering
    \includegraphics[width=0.49\textwidth]{state_horiz.png}
    \caption{(слева) Карта проезжей части аэропорта: зеленые ячейки (=1) — разрешенные ячейки,
			  темные клетки (=0) запрещены к перемещению, желтая клетка — пункт назначения на рейсе (=10);
			  (справа) положение рейса на карте в иде матрицы: желтая ячейка (=1) обозначает позицию самолета (=1),
			  а темная ячейка (=-1) описывает текущее направление движения}
    \label{fig:state_map}
\end{figure}

\subsection{Свойства и действия агентов}

Каждый рейс или Агент имеет позицию на карте в виде ячейки сетки, в которой он находится в дискретный момент $t$.

\begin{equation}
    \textbf{s}_t = \{s^i_t\}^N_{i=1}
\end{equation}

Мы считаем, что максимальное число агентов $N$ равно ограничено количеством ворот.
Чтобы обеспечить более свободное перемещение по карте, агенту
было добавлено новое свойство - базовое направление $i$-го полета.

\begin{equation}
    \textbf{v}_t = \{v^i_t\}^N_{i=1}, v \in (\{-1,0,1\} \times \{-1, 0, 1\})\backslash(0,0)
\end{equation}

Это необходимо для применения направленных движений по карте, делённых на три группы:

\begin{itemize}
\item Движения без изменения направления: шаг вперед в текущем направлении полета и остановка,
  фиксирующая положение самолёта на один временной шаг
\end{itemize}
\begin{equation}
    \begin{cases}
        s^i_{t+1} = s^i_{t} + v^i_t * f \\
        f = 1 \text{ if step = 'forward', else } 0 \\
        v^i_{t+1} = v^i_{t}
    \end{cases}
\end{equation}
\begin{itemize}
\item Небольшие повороты: движения с поворотом базового направления на 45 градусов и
  шагом вперед в измененном направлении
\end{itemize}
\begin{equation}
    \begin{cases}
        v^i_{t+1} = rotate(v^i_{t}, \pm 45^{\circ}) \\ 
        s^i_{t+1} = s^i_{t} + v^i_{t+1}\\
    \end{cases}
\end{equation}
\begin{itemize}
\item Крутые повороты: движения в текущем направлении, но с последующим изменением конечного направления 
 на 90 градусов.
\end{itemize}
\begin{equation}
    \begin{cases}
        s^i_{t+1} = s^i_{t} + v^i_{t}\\
        v^i_{t+1} = rotate(v^i_{t}, \pm 90^{\circ}) \\
    \end{cases}    
\end{equation}

Итоговый набор движений, возможный для совершения самолетом в одном направлении представлен на графике \ref{fig:moves}.
Так, будучи в точке $s_0 = (0,0)$ в направлении $v = (1,0)$, самолёт может продвинуться впёред, до точки (1,0) или, 
совершив небольшой поворот на 45 градусов, сместиться в точку (1,-1) или (1,1) - тогда его направление изменится
на 45 градусов по часовой или против часовой стрелки соответственно.
Наконец, для совершения более крутого поворота, самолёт-агент может продвинуться в точку (1,0) и сменить направление
до $(0,\pm 1)$. 

Применяемый комплекс действий подчеркивает инертный характер маневров самолета -
требуется несколько шагов и дополнительное пространство для разворота самолета
при движении по проезжей части.
Например, требуется выполнить два одинаковых резких поворота или четыре
небольших поворота, чтобы изменить направление самолета на противоположное.
Однако, поскольку в данной работе скорость самолета не учитывается, его остановка происходит всегда
мгновенно, но сопровождается штрафом $\gamma$.

Чтобы стимулировать скорейшее движение к выходам,
в окружение втроен фоновый штраф $\alpha$, который применяется к полету на каждом шаге.

Последним дополнением к свойствам полета является пункт назначения – ячейка
одного из выходов на карте проезжей части - эта точка является главной целью задания в окружении для каждого самолёта.

\begin{equation}
	d^i \in \{ (9,2), (9,5) \} = D
\end{equation}

Назначаемый для каждого самолёта выход,
который, как правило, располагается ближе всего к воротам, из которых появился рейс.
За достижение пункта назначения к успешным самолетам применяется положительная награда $\theta$.
Если самолет попадает к другому выходу, не запланированному для него, награда также выдаётся, но уменьшенная.


\begin{figure}[h]
\centering
\label{fig:moves}
\includegraphics[width=0.25\textwidth]{moves.png}
\caption{Карта всех возможных действий с позиции (0,0) и направления (1,0)
		  (синяя точка) и будущие состояния (оранжевые точки — позиции и стрелки — направления),
		  включая «остановку» на том же месте.}
\end{figure}

\subsection{Системы наград и остановки эпизода}

Все вышеописанные ранее правила окружения можно описать в системе уравнений наград:

\begin{equation*}
r^i_t = -\alpha +
\begin{cases}
-\gamma, \text{ if } s^i_{t-1} = s^i_{t} \\
-\beta, \text{ if } \exists \tau_0: s^i_t + \tau_0 u^i_t = s^j_t + \tau_0 u^j_t\\
\ \ \ \ \ \ \ \ \ i \neq j, \tau_0 \in (0,1] \\
-\phi,  \text{ if } s^i_t \text{ is not a permitted cell }  \\
\ \ \theta, \text{ if } s^i_t = d^i \\
\ \ \theta_r, \text{ if } s^i_t \in D \backslash d^i \\
\end{cases}
\end{equation*}

\begin{equation}
    \alpha, \beta, \phi, \theta, \theta_r, \gamma > 0,\ \theta > \theta_r,\ \gamma > \alpha
\end{equation}

Если самолет получает одну из условных наград (любую кроме базовой награды $-\alpha$ и штрафа за остановку $-\gamma$),
 самолет становится неактивным, а его последующие награды аннулируются до конца эпизода.
Однако агенты могут покинуть карту окружения только через выходы справа, 
поэтому после его выхода из строя, даже в неактивном режиме он все равно останется препятствием для 
других активных агентов (ячейка, где остановился вышедший из строя самолет, становится неразрешенной).
Эпизод заканчивается после того, как все рейсы станут неактивными.
Продолжительность эпизода ограничена 12 шагами,
поэтому, если какой-либо из самолётов остался активным до "дедлайна", эпизод обрывается.